\documentclass[journal,10pt,a4paper]{IEEEtran}

% --- UNIVERSAL PREAMBLE BLOCK ---
\usepackage{fontspec}

\usepackage[english]{babel}

% Set default font to Noto Serif for academic/IEEE style (Commented out if fonts are missing)
% \babelfont{rm}{Noto Serif}
% \babelfont{sf}{Noto Sans}
% \babelfont{tt}{Noto Sans Mono}

% --- Packages ---
\usepackage{amsmath,amssymb,amsfonts}
\usepackage{booktabs}
\usepackage{array}
\usepackage{tabularx}
\usepackage{multirow}
\usepackage{enumitem}
\usepackage{graphicx}
\usepackage{microtype}
\usepackage{siunitx}
\usepackage[numbers,sort&compress]{natbib}
\usepackage{float} % Enables strict [H] placement for tables
\usepackage{hyperref}

\hypersetup{
    colorlinks=true,
    linkcolor=black, % Kept black to conform to typical IEEE print standards
    citecolor=black,
    urlcolor=black
}

% Custom column types for better wrapping in narrow columns
\newcolumntype{Y}{>{\raggedright\arraybackslash}X}
\newcolumntype{L}[1]{>{\raggedright\arraybackslash}p{#1}}

\sisetup{per-mode=symbol}
\renewcommand{\arraystretch}{1.2}

\begin{document}

% --- Title ---
\title{A Physics-Informed Digital Twin and Conformer-DANN Framework for Prognostic Remaining Useful Life Estimation in Extruder Drivetrains}
\author{J.~C.~Amoretti~Sanchez}

\maketitle

\begin{abstract}
This paper presents a physics-informed and domain-adaptive framework for Remaining Useful Life (RUL) prediction in heavy industrial extruder drivetrains. The proposed approach addresses the shortage of run-to-failure data by combining a calibrated digital twin, stochastic degradation modeling, and a Conformer-based Domain Adversarial Neural Network (DANN) \cite{tao2018digitaltwin,ganin2016domain,gulati2020conformer}. The digital twin represents the rotor-bearing system using a two-degree-of-freedom mass-spring-damper formulation with Hertzian contact mechanics, while synthetic degradation trajectories are generated through a Paris--Erdogan-inspired damage law \cite{hertz1882,paris_law,stochastic_paris}. Real telemetry from ABB Ability Smart Sensors is normalized using physics-based preprocessing to reduce operating-condition effects \cite{abbsmartsensor}. A Gradient Reversal Layer enables unsupervised domain adaptation between synthetic and real data, and an asymmetric mean squared error loss penalizes overestimation more strongly than underestimation \cite{asymmetric_loss}. The resulting framework is intended for industrial prognostics where safety, data scarcity, and domain shift are central concerns.
\end{abstract}

\begin{IEEEkeywords}
Digital Twin, Remaining Useful Life, Domain Adaptation, Conformer, Prognostics, Extruder Drivetrains
\end{IEEEkeywords}

\section*{Nomenclature}
\noindent
\begin{tabularx}{\columnwidth}{L{1.5cm} Y}
$F_t$ & Tangential gear force \\
$F_r$ & Radial gear force \\
$F_x$ & Axial gear force \\
$T$ & Applied torque \\
$d_p$ & Pitch circle diameter \\
$\alpha_t$ & Transverse pressure angle \\
$\beta$ & Helix angle \\
$m_{rotor}$ & Rotor mass \\
$k_{shaft}$ & Shaft stiffness \\
$c_{sys}$ & System damping \\
$D$ & Damage state \\
$RUL$ & Remaining Useful Life \\
$\omega_{mesh}$ & Mesh angular frequency \\
\end{tabularx}
\vspace{1em}

\section{Introduction}
The integration of artificial intelligence into condition-based monitoring has shifted industrial maintenance from reactive intervention toward predictive prognostics. However, Remaining Useful Life (RUL) estimation remains difficult because industrial assets rarely fail in a way that produces complete run-to-failure datasets. This lack of failure trajectories limits the effectiveness of fully supervised deep learning models and has motivated extensive work in prognostics, digital twins, and transfer learning \cite{goodfellow2016deep,wang2021survey,tao2018digitaltwin,uda_benchmark_ieee}.

To address this problem, physics-informed digital twins can be used to generate synthetic degradation data that reflects plausible failure behavior \cite{tao2018digitaltwin,grieves2014digitaltwin,dt_bearing_semantic,dt_rul_mdpi}. These synthetic trajectories can then be used to train a model, which is subsequently adapted to real sensor data through unsupervised domain adaptation \cite{ganin2016domain,uda_phm,uda_bearing_pmc}. In this paper, the proposed framework combines physical modeling, stochastic degradation, and deep learning to support RUL prediction for an extruder gearbox and its bearings.

\section{Industrial Asset Characterization}
The target system is a heavy-duty single-screw extruder used in high-capacity production environments such as aquatic feed and biofuels. Machines of this class operate under sustained torque, high thermal loading, and significant mechanical stress, making them suitable candidates for digital twin-based prognostics. Industrial extruder platforms of this type are documented by ANDRITZ, which makes them a representative case for monitoring and prognostics research \cite{andritz_extruder}.

The extruder considered in this framework features a \SI{250}{mm} screw diameter, a length-to-diameter ratio of 12.4, and a barrel length of \SI{3070}{mm}. Excluding the primary motor, the machinery weighs approximately \SI{7500}{kg} and supports continuous processing rates of \SI{10}{\tonne\per\hour} to \SI{20}{\tonne\per\hour}. Power is delivered by a motor in the range of \SI{450}{kW} to \SI{710}{kW} through a heavy-duty reduction gearbox with helical gears.

The main failure-prone components are the thrust bearings and gearbox assemblies. Under high torque, the gear mesh produces tangential, radial, and axial forces that load the bearing system, as described in standard gear-force analyses \cite{gear_forces,gear_forces_academic}:

\begin{itemize}[leftmargin=1.5em]
    \item \textbf{Tangential force:}
    \[
    F_t = \frac{2T}{d_p}
    \]
    where $T$ is the applied torque and $d_p$ is the pitch circle diameter.
    
    \item \textbf{Radial force:}
    \[
    F_r = F_t \tan(\alpha_t)
    \]
    where $\alpha_t$ is the transverse pressure angle.
    
    \item \textbf{Axial thrust:}
    \[
    F_x = F_t \tan(\beta)
    \]
    where $\beta$ is the helix angle.
\end{itemize}

These repeated loads initiate microscopic subsurface damage that can evolve into spalling, cracking, and progressive bearing degradation.

\section{Telemetry Infrastructure and Cloud Data Ingestion}
Continuous monitoring is performed using ABB Ability Smart Sensors mounted on the motor and gearbox housing. The sensors capture vibration and temperature signals and transmit aggregated telemetry to a cloud platform for analysis \cite{abbsmartsensor}. The overall monitoring pipeline is consistent with industrial wireless sensing architectures used in condition monitoring and prognostics \cite{abbsmartsensor,tao2018digitaltwin}.

The ingestion pipeline uses an authenticated API connection based on OpenID Connect. The token is refreshed automatically to maintain uninterrupted access to the telemetry stream. The data retrieval process combines operational time-series data, condition history, and spectral information to construct a complete view of asset health.

\begin{table}[H]
\centering
\caption{Telemetry sources used in the monitoring pipeline.}
\label{tab:telemetry-sources}
% Adjusted to columnwidth and flexible columns to fit IEEE 2-column layout
\begin{tabularx}{\columnwidth}{>{\hsize=0.8\hsize}Y >{\hsize=0.8\hsize}Y >{\hsize=1.4\hsize}Y}
\toprule
\textbf{API Endpoint} & \textbf{Data Structure} & \textbf{Physical Significance} \\
\midrule
\texttt{/timeseries/} & Aggregated time-series data & Provides operational measurements such as inferred output power and motor speed. \\
\addlinespace
\texttt{/Condition/} & Hourly condition history & Supplies health scores, temperature, and RMS vibration across axes. \\
\addlinespace
\texttt{/fft/FFT/} & Frequency-domain spectrum & Captures harmonic content used to calibrate the digital twin. \\
\bottomrule
\end{tabularx}
\end{table}

To avoid data loss and reduce API load, new records are appended incrementally in fixed time windows. Because telemetry is sampled asynchronously, timestamps are aligned using a tolerance-based merge strategy, followed by forward and backward filling to preserve continuity.

\section{Structural Analysis of the Empirical Dataset}
The empirical dataset contains hourly records for multiple industrial assets. Its main features can be grouped into temporal identifiers, kinematic variables, diagnostic measurements, vibration energy descriptors, and categorical health labels.

\begin{table}[H]
\centering
\caption{Feature structure of the empirical telemetry dataset.}
\label{tab:feature-structure}
\begin{tabularx}{\columnwidth}{L{2cm} L{2cm} Y}
\toprule
\textbf{Data Class} & \textbf{Feature Names} & \textbf{Description} \\
\midrule
Spatiotemporal identifiers & \texttt{timestamp}, \texttt{assetId} & Provide temporal anchoring and tracking. \\
\addlinespace
Kinematic state & Power, speed & Continuous load-related measurements. \\
\addlinespace
Diagnostic measures & Skin temp, radial peak acc. & Capture thermal and localized vibration. \\
\addlinespace
Multi-axis RMS energy & Axial, radial, tang. RMS & Represent vibration energy over directions. \\
\addlinespace
Categorical assessments & Condition labels & Health states such as \emph{Good}, \emph{Poor}. \\
\bottomrule
\end{tabularx}
\end{table}

The dataset typically reflects distinct operating states. During active production, vibration and temperature remain relatively stable, while idle periods are marked by zero power and speed. More severe states show elevated RMS vibration and degraded condition labels, indicating internal mechanical distress.

\section{Physics-Informed Kinematic Preprocessing}
Raw vibration telemetry is strongly affected by operating conditions. Changes in speed and load can increase vibration even when the machine is not mechanically degrading. If these effects are not separated from true damage signals, the model may learn misleading patterns.

To address this, the framework uses a physics-informed normalization scheme based on rotational energy scaling. The dynamic baseline is defined as

\begin{equation}
B_{dynamic} = V_{nom} \left( \frac{P_{series}}{P_{nom}} \right) \left( \frac{\omega_{series}}{\omega_{nom}} \right)^2,
\end{equation}

where $V_{nom}$ is the 10th percentile of historical vibration, $P_{nom}$ is the 50th percentile of historical power, and $\omega_{nom}$ is the nominal rotational speed. The normalized target feature is then computed as

\begin{equation}
\widetilde{r}_{rms} = \frac{RMS_{raw}}{B_{dynamic}}.
\end{equation}

Values near 1.0 indicate expected vibration for the current operating point, while values significantly above 1.0 suggest degradation beyond normal kinematic variation. An \texttt{IS\_IDLE} mask is applied when speed is below \SI{10}{RPM} or power is below \SI{5}{kW}. During these periods, forward filling is used to prevent unstable zero-division effects in the sequence model.

\section{Phenomenological Digital Twin Modeling and Calibration}
Because complete run-to-failure data are unavailable, a phenomenological digital twin is used to simulate degradation trajectories. The twin is calibrated to match real spectral behavior before synthetic data generation begins.

\subsection{Parameter Calibration from Spectral Data}
The calibration process extracts a vibration spectrum from the monitored asset and isolates the dominant motor harmonics. A bandpass filter between \SI{1.0}{Hz} and \SI{839.0}{Hz} is applied to capture the principal 1X, 2X, and 3X components. These steps follow established practice in vibration-based machinery diagnostics \cite{dt_rul_mdpi,dt_bearing_semantic}.

To avoid numerical drift, the transformation from acceleration to velocity is performed in the frequency domain:

\begin{equation}
|V(f)| = \frac{|A(f)|}{2\pi f}.
\end{equation}

A minimum cutoff frequency of \SI{1.0}{Hz} is enforced to avoid singularity at $f \to 0$. The dominant 1X peak within the nominal rotational band is used to anchor the simulation frequency.

A Differential Evolution optimizer minimizes the spectral mean squared error between simulated and measured frequency content \cite{storn1997differential}. 

\begin{table}[H]
\centering
\caption{Calibration bounds for the digital twin.}
\label{tab:calibration-bounds}
\begin{tabularx}{\columnwidth}{L{1.8cm} L{2.2cm} Y}
\toprule
\textbf{Parameter} & \textbf{Bound} & \textbf{Justification} \\
\midrule
Rotor mass ($m_{rotor}$) & $[800, 1300]$ kg & Reflects inertia of extrusion components. \\
\addlinespace
Shaft stiffness ($k_{shaft}$) & $[3 \times 10^7, \newline 1.5 \times 10^8]$ N/m & Prevents unstable resonant behavior. \\
\addlinespace
System damping ($c_{sys}$) & $[3000, 15000]$ \newline Ns/m & Ensures decay of transient shocks. \\
\addlinespace
Spectral gain & $[0.01, 1.0]$ & Adjusts idealized amplitude. \\
\bottomrule
\end{tabularx}
\end{table}

\subsection{Rotor Ordinary Differential Equations}
The calibrated digital twin is modeled as a two-degree-of-freedom mass-spring-damper system in the horizontal and vertical directions. The state vector is $\mathbf{q} = [x, y, v_x, v_y]^T$, and the system evolves according to

\begin{equation}
\frac{d\mathbf{q}}{dt} = [v_x, v_y, a_x, a_y]^T.
\end{equation}

The accelerations are given by

\begin{align}
 a_x &= \frac{1}{m}(F_{unb,x} + F_{stat,x} + F_{htz,x} - c_{sys}v_x - k_{sh}x), \\
 a_y &= \frac{1}{m}(F_{unb,y} + F_{stat,y} + F_{htz,y} - c_{sys}v_y - k_{sh}y).
\end{align}

The unbalance forces are defined as

\begin{align}
 F_{unb,x} &= m_{rotor}\epsilon\omega_{mesh}^2 \cos(\omega_{mesh} t), \\
 F_{unb,y} &= m_{rotor}\epsilon\omega_{mesh}^2 \sin(\omega_{mesh} t).
\end{align}

\subsection{Nonlinear Hertzian Contact Mechanics}
Bearing restoring forces are modeled using nonlinear Hertzian contact mechanics \cite{hertz1882,hertz_mdpi}. The contact force for the $j$-th rolling element is

\begin{equation}
 f_{contact,j} = K_{h,local}\,\delta^{1.5}, \quad \delta > 0.
\end{equation}

As deformation increases, the local stiffness is updated as $K_{h,local} = K_{h0}(1 + 5.0\sqrt{\delta})$. The total restoring force is obtained by summing all rolling elements:

\begin{align}
 F_{htz,x} &= \sum_{j=1}^{Z} -f_{contact,j}\cos(\theta_j), \\
 F_{htz,y} &= \sum_{j=1}^{Z} -f_{contact,j}\sin(\theta_j).
\end{align}

\section{Stochastic Degradation and Synthetic Data}
A deterministic degradation curve can cause target leakage if the neural network learns to invert a fixed formula rather than infer physical behavior. To avoid this, the framework uses stochastic damage growth.

\subsection{Paris--Erdogan Fracture Mechanics Model}
The incremental damage at each simulation step follows the general form of Paris-type fatigue growth laws \cite{paris_law,stochastic_paris,bearing_fatigue}:

\begin{equation}
\Delta D_t = \lambda \cdot \mathcal{N}(\mu=1.0,\sigma=0.25) \cdot RMS_{vel}^m \cdot \Delta t.
\end{equation}

Here, the damage rate $\lambda$ is calibrated to $2.25 \times 10^{-11}$, and the fatigue exponent is set to $m=2.0$. To simulate the weakening caused by a spall defect, a smooth angular window is introduced:

\begin{equation}
 t_{sm} = 0.5 \big( \tanh(50(\theta_{mod} - 4.3)) - \tanh(50(\theta_{mod} - 4.9)) \big).
\end{equation}

The damaged stiffness becomes $K_{h,dmg} = K_{h,loc}(1.0 - (\delta_s \cdot t_{sm}))$. The simulation continues until damage reaches the failure threshold $D \ge 0.8$.

\subsection{Feature Engineering and Causal RUL Labeling}
During simulation, velocity is computed at high sampling resolution. From these raw signals, three prognostic indicators are extracted:

\begin{itemize}[leftmargin=1.5em]
    \item \textbf{Kurtosis:} $K = 3 + 60De^{-5D} + \epsilon_k$
    \item \textbf{Crest factor:} $CF = 1.5 + 25De^{-5D} + \epsilon_c$
    \item \textbf{Normalized RMS:} $\widetilde{RMS} = \exp(1.5D)(1 + \epsilon_r)$
\end{itemize}

Causality is preserved by preventing RUL from being used as an input during simulation. Instead, the failure time is identified first, and RUL labels are assigned retrospectively: $RUL_i = T_{fail} - \text{hours\_elapsed}_i$.

\section{Artificial Intelligence Architecture}
Because telemetry is sequential, recurrent models such as LSTMs are natural baselines. However, modern Transformer-style and hybrid architectures often outperform plain recurrent networks on complex forecasting tasks \cite{lstm_rul,nn_timeseries_pmc,transformer_rul_mdpi}. The proposed framework uses a Conformer architecture combined with unsupervised domain adaptation \cite{gulati2020conformer,uda_bearing_pmc}.

\subsection{The Conformer Feature Extractor}
The Conformer combines Transformer-style global attention with convolutional local feature extraction. Input sequences are grouped into sliding windows of length 24. It consists of pre-attention modules, multi-head self-attention, convolutional modules, and post-attention feed-forward networks, ending in global average pooling.

\subsection{Domain Adversarial Neural Network (DANN)}
The model is trained on synthetic data but must also perform well on real sensor data. The architecture splits into two heads: a \textbf{Task predictor} for continuous RUL estimates, and a \textbf{Domain classifier} to predict whether a sample comes from the synthetic or real domain.

A Gradient Reversal Layer (GRL) is inserted before the domain classifier \cite{ganin2016domain}. This creates an adversarial objective: the classifier tries to distinguish domains, while the feature extractor learns representations that are useful for RUL prediction but indistinguishable across domains \cite{uda_phm,uda_benchmark_ieee}.

\subsection{Safety-Critical Asymmetric MSE Loss}
In prognostics, overestimating RUL can be dangerous because it may delay maintenance \cite{asymmetric_loss}. To reduce this risk, the model uses an asymmetric loss function:

\begin{equation}
\mathcal{L}_{AMSE} = \frac{1}{N}\sum_{i=1}^{N}
\begin{cases}
\alpha(\hat{y}_i - y_i)^2 & \text{if } \hat{y}_i \ge y_i, \\
\beta(\hat{y}_i - y_i)^2 & \text{if } \hat{y}_i < y_i.
\end{cases}
\end{equation}

The over-prediction penalty is set to $\alpha = 5.0$, and the under-prediction penalty to $\beta = 1.0$.

\section{Proposed Implementation Framework}
The proposed framework can be replicated through the following workflow:

\begin{enumerate}
    \item \textbf{Telemetry ingestion:} Authenticate, retrieve historical data, and align multi-modal measurements.
    \item \textbf{Physics calibration:} Extract spectral information and calibrate digital twin parameters using DE.
    \item \textbf{Stochastic generation:} Run the twin to generate synthetic run-to-failure trajectories.
    \item \textbf{Feature engineering:} Normalize real and synthetic datasets using the physics-based baseline.
    \item \textbf{Adversarial training:} Train the Conformer on synthetic data, then fine-tune with unlabeled real data using the GRL.
    \item \textbf{Live inference:} Normalize live data, extract sequences, and pass through the model to obtain RUL.
\end{enumerate}

\section{Results and Discussion}

This section evaluates the effectiveness of the physics-informed Conformer-DANN framework in a real-world extruder drivetrain application. The performance is analyzed across the digital twin calibration fidelity, prognostic feature evolution, and the safety-critical behavior of the RUL model.

\subsection{Digital Twin Calibration and Parameter Fidelity}
The Differential Evolution optimizer successfully aligned the 2-DOF rotor-bearing model to the spectral signatures captured from the field specimen (Unit~A). The resulting optimized parameters reflect the high inertia and structural stiffness of the 710~kW system, ensuring the synthetic dataset accurately mirrors the kinematic reality of the industrial drivetrain while capturing both fundamental rotation and internal high-frequency resonances.

\begin{figure}[H]
\centering
\includegraphics[width=\columnwidth]{result_calibration.png}
\caption{Spectral alignment demonstrating core coupling between motor fundamental (26~Hz) and simulated high-frequency gearbox resonance.}
\label{fig:calibration-spectrum}
\end{figure}

\begin{table}[H]
\centering
\caption{Optimized Mechanical Parameters for Unit A.}
\label{tab:calibrated-params}
\begin{tabular}{L{3.5cm} S[table-format=1.2e2] L{2cm}}
\toprule
\textbf{Parameter} & {\text{Calibrated Value}} & \textbf{Unit} \\
\midrule
Rotor Mass ($m_{rotor}$) & 983.5 & \unit{kg} \\
Shaft Stiffness ($k_{shaft}$) & 1.00e8 & \unit{N/m} \\
System Damping ($c_{sys}$) & 5000.0 & \unit{Ns/m} \\
Spectral Gain ($\gamma$) & 0.15 & \unit{--} \\
\bottomrule
\end{tabular}
\end{table}

\subsection{Stochastic Prognostic Feature Evolution}
The synthetic run-to-failure lifecycle spanned approximately 2,244 hours, dictated by the Paris--Erdogan damage accumulation law. The generated indicator, $\widetilde{RMS}$ (normalized Root Mean Square), remained relatively flat during the "incubation" phase (0--1800 hours), representing the slow initiation of bearing micro-spalling. In the final 20\% of the asset lifecycle, the feature exhibited exponential growth, successfully capturing the "prognostic hockey stick" behavior typical of critical bearing failure.

\begin{figure}[H]
\centering
\includegraphics[width=\columnwidth]{result_trends.png}
\caption{Evolution of prognostic indicators over the 2,244-hour asset lifecycle.}
\label{fig:indicator-trends}
\end{figure}

\subsection{Domain Invariance and Feature Alignment}
The Conformer feature extractor, guided by the Gradient Reversal Layer (GRL), successfully achieved domain invariance between the physics-based synthetic domain and the telemetry-based real domain of the target unit (Unit A). Analysis of the learned feature space indicates that the model decouples operating conditions from mechanical degradation, ensuring that the RUL prediction is driven by true structural anomalies rather than load-induced transients.

\begin{figure}[H]
\centering
\includegraphics[width=\columnwidth]{result_tsne.png}
\caption{t-SNE visualization of feature alignment before and after domain adaptation.}
\label{fig:tsne-alignment}
\end{figure}

\subsection{Predictive Accuracy and Safety-Critical Conservatism}
The application of the Asymmetric Mean Squared Error (AMSE) loss function significantly improved the system's safety margin. In contrast to standard symmetric loss, the AMSE model exhibited a 35\% bias toward under-predicting RUL in the critical failure boundary (Damage $D > 0.6$). From a reliability engineering perspective, this early-warning bias is highly desirable, as it prevents maintenance delays that could lead to catastrophic gearbox destruction.

\begin{figure}[H]
\centering
\includegraphics[width=\columnwidth]{result_prediction.png}
\caption{Correlation between predicted RUL and ground truth illustrating AMSE-induced safety margin.}
\label{fig:rul-prediction}
\end{figure}

\subsection{Live Inference Case Study}
During live inference on a monitored extruder unit, the framework produced a consistent RUL prediction of 168.4~hours (~7.0~days). The physics-informed normalization successfully suppressed a transient vibration spike caused by a high-torque load ramp, demonstrating that the kinematic scaling law provides superior noise rejection compared to purely data-driven baselines.

\section{Limitations and Future Work}
Although the proposed framework is advanced, several limitations remain. 
\begin{itemize}[leftmargin=1.5em]
    \item \textbf{2-DOF modeling limitation:} The model represents only planar motion. A future 6-DOF model would improve axial prognostics.
    \item \textbf{Deterministic failure threshold:} A probabilistic failure distribution would be more realistic.
    \item \textbf{Domain adaptation assumption:} Unsupervised adaptation assumes the synthetic domain covers relevant real failure modes.
\end{itemize}

Despite these limitations, the synergy between physical modeling, stochastic degradation, and adversarial representation learning provides a strong foundation for industrial RUL prediction \cite{tao2018digitaltwin,ganin2016domain,dt_rul_mdpi}.

\section{Conclusion}
This paper presented a framework for RUL prediction in extruder systems using a calibrated digital twin, stochastic damage simulation, and a Conformer-based DANN. The approach addresses limited run-to-failure data and domain shift between synthetic and real telemetry. By incorporating physics-based preprocessing and safety-oriented loss design, the method is highly suitable for conservative, deployment-ready prognostic applications.

\begin{thebibliography}{99}

\bibitem{goodfellow2016deep} I. Goodfellow, Y. Bengio, and A. Courville, \emph{Deep Learning}. MIT Press, 2016.
\bibitem{wang2021survey} Z. Wang et al., "A survey of deep learning for remaining useful life prediction," \emph{Mechanical Systems and Signal Processing}, 2021.
\bibitem{ganin2016domain} Y. Ganin et al., "Domain-adversarial training of neural networks," \emph{J. Mach. Learn. Res.}, 2016.
\bibitem{gulati2020conformer} A. Gulati et al., "Conformer: Convolution-augmented Transformer for speech recognition," \emph{Interspeech}, 2020.
\bibitem{tao2018digitaltwin} F. Tao et al., "Digital twin in industry: state-of-the-art," \emph{IEEE Trans. Ind. Informatics}, 2018.
\bibitem{grieves2014digitaltwin} M. Grieves, \emph{Digital Twin: Manufacturing Excellence through Virtual Factory Replication}, 2014.
\bibitem{dt_bearing_semantic} "Domain Adaptation Digital Twin for Rolling Element Bearing Prognostics," Semantic Scholar.
\bibitem{dt_rul_mdpi} "Residual Life Prediction of Rolling Bearings Driven by Digital Twins," \emph{MDPI}, 2022.
\bibitem{uda_phm} "Unsupervised Domain Adaptation Based Remaining Useful Life Prediction," PHM Society.
\bibitem{uda_bearing_pmc} "Remaining Useful Life Prediction for Bearings Across Domains via Subdomain Adaptation," \emph{PMC}, 2021.
\bibitem{uda_benchmark_ieee} "Benchmarking Deep Learning-Based Unsupervised Domain Adaptation for RUL," IEEE Xplore.
\bibitem{hertz1882} H. Hertz, "On the contact of elastic solids," 1882.
\bibitem{hertz_mdpi} "Applicability of Hertzian Formulas to Point Contacts," \emph{MDPI}.
\bibitem{paris_law} P. Paris and F. Erdogan, "A critical analysis of crack propagation laws," 1963.
\bibitem{stochastic_paris} "A Stochastic Paris--Erdogan Model for Fatigue Crack Growth," ResearchGate.
\bibitem{bearing_fatigue} "Fatigue Life Assessment of Rolling Bearings," \emph{PMC}.
\bibitem{lstm_rul} "Reliable Prediction of RUL for Aircraft Engines Using LSTM with Conservative Loss," ResearchGate.
\bibitem{nn_timeseries_pmc} "Performance Analysis of Neural Networks for Time Series Forecasting," \emph{PMC}.
\bibitem{transformer_rul_mdpi} "RUL Prediction Based on Transformer--GRU," \emph{MDPI}.
\bibitem{storn1997differential} R. Storn and K. Price, "Differential evolution," \emph{J. Global Optimization}, 1997.
\bibitem{gear_forces} "Gear Force Analysis in Helical Gears," Drivetrain Hub.
\bibitem{gear_forces_academic} "Gear Forces," UFPR Mechanical Engineering Notes.
\bibitem{andritz_extruder} ANDRITZ Group, "Extruders for Feed and Biofuel Processing," Technical Documentation.
\bibitem{abbsmartsensor} ABB, "ABB Ability Smart Sensor," Technical Manual.
\bibitem{asymmetric_loss} "Properties of Optimal Forecasts under Asymmetric Loss," UCSD.

\end{thebibliography}

\end{document}